% !PRINTED 293
\weberpage{293}
%<<CONT
sind auch alle folgenden Gruppen $P_{\lambda+1}, P_{\lambda+2}, \ldots$, die ja alle
Theiler von $P_\lambda$ sind, intransitiv.

Es ist nun an die Sätze 1., 2., 3. im §.\,158 des ersten Bandes zu erinnern. Da
$P_\lambda$ ein Normaltheiler von $P_{\lambda-1}$ ist, so folgt aus jenen Sätzen,
dass $P_{\lambda-1}$ imprimitiv ist, und wenn durch die nöthigen Adjunctionen die
Gruppe von $f(x) = 0$ auf $P_{\lambda-1}$ reducirt ist, so wird durch weitere
Adjunction einer Wurzel einer Normalgleichung vom Grade $j_\lambda$ die Function
$f(x)$ in mehrere irreducible Factoren
\weq{2}{f(x) = f_1(x)\, f_2(x) \ldots f_h(x)}
zerfallen, die alle von gleichem Grade sind und deren Anzahl $h$ ein Theiler von
$j_\lambda$ ist.

Bezeichnen wir mit $m$ den Grad von $f(x)$, mit $m_1$ den Grad von $f_1(x)$, so ist
\weq{3}{m = h\, m_1.}
Nach dem angeführten Satze 1. im §.\,158, Bd. I haben die Gleichungen
$f_1 = 0, f_2 = 0, \ldots, f_h = 0$ alle dieselbe Gruppe, die wir mit $Q$ bezeichnen
wollen. Sind $\alpha, \alpha_1, \alpha_2, \ldots, \alpha_{m_1-1}$ die Wurzeln von
$f_1 = 0$, so erhalten wir die Gruppe $Q$, wenn wir die Permutationen der $\alpha$
sammeln, die durch $P_\lambda$ hervorgerufen werden. Es handelt sich zunächst um
das Verhältniss der Grade $p_\lambda$ und $q$ der beiden Gruppen $P_\lambda$ und $Q$.

Es kann mehrere verschiedene Permutationen in $P_\lambda$ geben, die dasselbe
Element in $Q$ erzeugen, die also dieselbe Permutation der $\alpha$ enthalten.
Sind $\pi_1, \pi_2$ zwei verschiedene Permutationen aus $P_\lambda$, die unter den
$\alpha$ dieselbe Permutation hervorrufen, so wird $\pi_2 \pi_1^{-1} = \pi_0$ eine
Permutation sein, die die $\alpha$ in Ruhe lässt, und es ist also
$\pi_2 = \pi_0 \pi_1$. Wenn umgekehrt $\pi_0$ irgend eine Permutation aus
$P_\lambda$ ist, die die Wurzeln $\alpha$ nicht permutirt, so wird die Permutation
$\pi_2 = \pi_0 \pi_1$ dieselbe Aenderung unter den $\alpha$ bedingen, wie $\pi_1$.

Es folgt hieraus, dass jede Permutation der $\alpha$ gleich oft durch die
Permutationen von $P_\lambda$ erzeugt wird, nämlich ebenso oft, als die $\alpha$
durch Permutationen aus $P_\lambda$ ungeändert bleiben. Bezeichnen wir diese Zahl
mit $p_0$, so ist also
\weq{4}{p_\lambda = p_0\, q,}
\gesperrt{oder der Grad $p_\lambda$ der Gruppe $P_\lambda$ ist ein Vielfaches des Grades $q$
von $Q$.}
